\section{管辖}
\begin{enumerate}
  \item 法院之间怎么分工。四级法院：最高、高、中、基层。
  \item 级别管辖、地域管辖。同级法院之间的分工。
  \item \zd{管辖的前提是法院主管管辖的前提是法院主管管辖的前提是法院主管}。
  \item 管辖权恒定原则。\zd{以起诉时为准}。
        \begin{enumerate}
          \item 对于原告要调整。
                \begin{enumerate}
                  \item 原告增加、变更诉讼请求。
                  \item \zd{防止原告通过先保留再增加规避级别管辖}，钻法律空子。
                \end{enumerate}
          \item 对于被告不允许调整。
                \begin{enumerate}
                  \item 反诉是独立的诉。不认可管辖完全可以另行起诉。
                  \item 反诉是\zd{牵连管辖}，直接向受理本诉的法院提。
                \end{enumerate}
        \end{enumerate}
  \item 原则上基层法院。
        \begin{kk}{中院的管辖}
          \begin{enumerate}
            \item \zd{重大涉外}案件、本辖区重大影响。
            \item 最高院确定由中院管辖的案件。
                  \begin{enumerate}
                    \item 海事、海上纠纷 \yjt 海事法院。
                    \item 公益讼案件。
                    \item 专利纠纷 \yjt 知识产权法院（\zd{最高院确定}的中院 \hz 基层法院）。
                  \end{enumerate}
            \item \zd{看见仲裁找中院，除了国内仲裁的保全/证据保全 \yjt 基层法院}。
          \end{enumerate}
        \end{kk}

        \begin{kk}{高院的管辖}
          \begin{enumerate}
            \item 本辖区范围重大影响的。
          \end{enumerate}
        \end{kk}

        \begin{kk}{最高院}
          \begin{enumerate}
            \item 全国范围重大影响。
            \item 认为应当由本院审理的案件。宣告性的意义，实践中这个法条从来没用过。
            \item \zd{有且仅有最高院有权管辖认为应当有本院管辖的案件}。
          \end{enumerate}
        \end{kk}
  \item 地域管辖。一般原告就被告。
        \begin{enumerate}
          \item 一般地域管辖：\zd{以当事人的住所与法院的隶属关系}确定的管辖。
                \begin{itemize}
                  \item 双方都有特殊情况用原则。
                  \item 被告一方有特殊情况，被告就原告。限\zd{身份关系}，防止原告\zd{滥诉}。
                \end{itemize}
          \item 特殊地域管辖。
                \begin{enumerate}
                  \item 侵权纠纷 \yjt 侵权行为地（\zd{扩大解释，沾边即可}）和被告住所地。
                  \item 合同纠纷 \yjt \zd{被告住所地和合同履行地}。
                        \begin{kk}{合同的履行地}
                          \qd{好像都可以根据常识判断……}
                          \begin{enumerate}
                            \item 合同没有履行 \& 约定的履行地没在当事人一方住所地 \yjt 被告住所地。
                            \item 约定履行地的 \yjt 约定。\qd{不管实际履行地是否和约定的一致，有约从约}。
                            \item 找实际履行地：
                                  \begin{itemize}
                                    \item 给付货币的 \yjt 接收货币方为合同履行地。
                                    \item 不动产的 \yjt 不动产所在地。
                                    \item 其他标的的 \yjt 履行义务方住所地。\zd{\sjt 结合合同性质看原告诉讼请求}。
                                    \item 财产租赁、融资租赁 \yjt 租赁物使用地。
                                    \item 信息网络方式交付标的 \yjt 买受人住所地；其他方式交付 \yjt 收货地。
                                  \end{itemize}
                          \end{enumerate}
                        \end{kk}
                  \item 其他特殊的。
                        \begin{enumerate}
                          \item 公司诉讼 \yjt 公司住所地。
                          \item 保险合同\zd{看保险标的物} \yjt 被告住所地、保险标的物所在地：
                                \begin{itemize}
                                  \item 人身保险 \yjt +被保险人住所地。
                                  \item 运输工具 \yjt +登记注册地。
                                  \item 运输中货物 \yjt +运输目的地、保险事故发生地。
                                \end{itemize}
                          \item 票据纠纷 \yjt 票据支付地。
                          \item 运输合同 \yjt 始发地、目的地、被告住所地。
                          \item 运输侵权 \yjt 被告住所地、事故发生地、最先到达地（以人为本\zd{就近原则}，车辆船舶航空器）。
                        \end{enumerate}
                \end{enumerate}
        \end{enumerate}
  \item 专属管辖。
        \begin{enumerate}
          \item 不动产纠纷 \yjt 不动产所在地。\zd{不动产的物权纠纷、农村土地承包经营、房屋租赁、政策性房屋买卖、建设工程施工}。
          \item 港口作业 \yjt 港口所在地。
          \item 遗产继承 \yjt 被继承人死亡时住所地和主要遗产所在地。
        \end{enumerate}
  \item 协议管辖。
        \begin{enumerate}
          \item 只适用于财产纠纷、必须书面。
          \item 选择的法院要与争议\zd{有实际联系}：原告、被告，签订、履行、标的物所在地。
          \item 不得违背级别管辖和专属管辖的规定。
          \item 约定两个以上。
                \begin{enumerate}
                  \item 都有选一个。
                  \item 要\zd{明确}。起诉时能根据协议来确定管辖。否则视为\zd{无效}。如由守约方地。
                \end{enumerate}
        \end{enumerate}
  \item \zd{做题：专属管辖 \yjt 协议管辖 \yjt 法定管辖 \yjt 被告住所地 \yjt 合同履行地}。
  \item 共同管辖。最先立案的法院，后立案应当移送给先立案的法院。
  \item 裁定管辖。由\zd{法院通过裁定}的方式来确定具体案件的管辖权。
        \begin{enumerate}
          \item 移送管辖。错误立案的\zd{纠错}程序。
                \begin{enumerate}
                  \item 已经受理案件，发现自己没有管辖权，同时受移送法院有管辖权。
                  \item 纠错纠错\qd{单方行为}，\zd{不需要}移送法院\zd{同意}。
                  \item 纠错纠错\zd{一次}就好，\zd{不允许退回 \hz 再移送} \yjt 报上级法院指定管辖。
                \end{enumerate}
          \item 指定管辖。\zd{实践中解决所谓的集体回避的问题}。
                \begin{enumerate}
                  \item 有管辖权法院由于\zd{特殊原因}不能行使管辖权的。
                  \item \zd{移送管辖}中受移送法院认为自己没有管辖权的。
                  \item \zd{管辖权争议} \yjt 争议双方 \yjt 两方高院（两次协商程序）\yjt 最高院。
                \end{enumerate}
          \item 管辖\zd{权}转移。上下级法院之间，对级别管辖的调整和\zd{变通}。
                \begin{enumerate}
                  \item 上可以审下。搞不定、新类型等，下认为上审更合适。\zd{下级报上级同意 \hz 上级直接下裁定}。
                  \item 上移下。\zd{确有必要；开庭前报自己的上级法院批准}。\qd{无论如何逃不过高院的监督}。
                \end{enumerate}
          \item \zd{考移送管辖和管辖权转移的区别}。区分整成表格没有任何意义，抓\zd{本质}。其实就是揪概念。
        \end{enumerate}
  \item 管辖权异议救济。
        \begin{enumerate}
          \item 本案\zd{一审当事人}。\zd{第三人绝对不准提管辖权异议}。
          \item 提交答辩状期间提出。\zd{其实是给被告做选择}。
          \item 如果应诉，应诉管辖，\zd{违反级别管辖和专属管辖除外}。就\zd{一次机会}，后续都不能再提。
          \item \zd{不予受理、驳回起诉和管辖权异议}唯三可上诉的裁定。
          \item 裁定作出之前，原告撤诉。在\zd{准子撤裁定书中一并说明}就行。
        \end{enumerate}
  \item 集中管辖。
        \begin{enumerate}
          \item 互联网法院。节约成本。
                \begin{enumerate}
                  \item 北京、广州、杭州 3 个，相当于基层法院。互联网纠纷、合同纠纷。
                  \item 先正常判断管辖，确定后有互联网法院的进行替换即可。
                \end{enumerate}
          \item 破产法院。
                \begin{enumerate}
                  \item 债务人：进入破产程序的那个法人。
                  \item 一律由受理破产申请的法院（破产法人住所地）集中管辖。
                  \item \zd{破产程序前已经受理}的无需移送，裁定诉讼中止。
                \end{enumerate}
        \end{enumerate}
  \item 担保合同中的管辖问题。一起起诉，根据\zd{主合同}确定。
\end{enumerate}
